\unchapter{Introduction}

\textbf{Topicality of the problem:}
In the context of growing digital transformation, Project Management Offices (PMOs) 
in Latvian enterprises increasingly rely on multiple information systems for planning, 
reporting, and governance. However, many organizations struggle with fragmented system 
landscapes where data must be manually consolidated across tools, leading to 
inefficiencies, data quality issues, and delayed decision-making. Despite the practical 
significance of this challenge, empirical research on information system integration 
maturity in PMO environments within Latvia remains limited.

\textbf{Aim of the study:}
The aim of the study is to assess the integration maturity of information systems 
within PMO environments in medium and large Latvian enterprises and examine its 
relationship with reporting automation, data quality, and governance standardization.

\textbf{Object of the study:}
Information systems within Project Management Office environments in Latvian enterprises.

\textbf{Subject of the study:}
The integration maturity of information systems and its relationship with reporting 
automation, data accuracy, and governance standardization in PMO environments.

\textbf{Research problem:}
Many PMOs in Latvian enterprises operate with low levels of information system 
integration, resulting in poor architectural coherence, manual data consolidation, 
reduced reporting efficiency, and inconsistent governance practices.

\textbf{Tasks of the study:}
\begin{enumerate}
    \item Conduct a systematic literature review on PMO digitalization, information 
    system integration architectures, and maturity models relevant to IS environments.
    \item Adapt and operationalize an Integration Maturity Model (Levels 1--4) with 
    clearly defined, survey-measurable criteria for each level, based on integration 
    mechanisms, automation level, reporting structure, and BI centralization.
    \item Design and validate a structured survey instrument containing Likert-scale 
    and categorical questions aligned with the Integration Maturity Model dimensions.
    \item Collect empirical data through distribution of the survey to representatives 
    of medium and large Latvian enterprises (50+ employees) with formal PMO structures, 
    targeting 30--50 responses.
    \item Conduct 2--3 structured expert interviews to supplement survey findings with 
    architectural depth and qualitative insights.
    \item Perform descriptive statistical analysis to determine the distribution of 
    integration maturity levels across the sampled PMO environments.
    \item Apply correlation analysis (Pearson or Spearman, depending on data 
    distribution) and simple regression analysis to examine relationships between 
    integration maturity and reporting automation, data accuracy, and governance 
    standardization.
    \item Interpret findings, draw conclusions regarding the hypotheses, and formulate 
    practical recommendations for PMO information system integration improvement in 
    Latvian enterprises.
\end{enumerate}

The correspondence between the research tasks and the thesis chapters is presented 
in Table~\ref{tab:task-chapter}.

{\singlespacing
\begin{table}[h]
  \caption{Correspondence of research tasks to thesis chapters}
  \label{tab:task-chapter}
  \centering
  \begin{tabular}{p{0.6\linewidth}p{0.3\linewidth}}
    \toprule
    \textbf{Research Task} & \textbf{Thesis Chapter}\\
    \midrule
    Task 1 -- Literature review on PMO digitalization and maturity models 
    & Chapter 1\\
    Task 2 -- Design of Integration Maturity Model (Levels 1--4) 
    & Chapter 1\\
    Task 3 -- Survey instrument design and validation 
    & Chapter 2\\
    Task 4 -- Empirical data collection from Latvian enterprises 
    & Chapter 3\\
    Task 5 -- Expert interviews 
    & Chapter 3\\
    Task 6 -- Descriptive statistical analysis 
    & Chapter 3\\
    Task 7 -- Correlation and regression analysis 
    & Chapter 3\\
    Task 8 -- Conclusions and practical recommendations 
    & Chapter 4\\
    \bottomrule
  \end{tabular}
\end{table}
}

\textbf{Hypotheses:}
\begin{enumerate}
    \item Most PMO environments in medium and large Latvian enterprises operate at 
    intermediate (Levels 2--3) rather than advanced integration maturity levels.
    \item PMOs utilizing API-based or platform-level integration demonstrate higher 
    reporting automation than those relying primarily on manual data consolidation.
    \item Integration maturity positively correlates with perceived data accuracy 
    and timeliness.
    \item The use of centralized BI/dashboard systems is associated with higher 
    governance standardization in PMO processes.
\end{enumerate}

\textbf{Research methods:}
The study applies systematic literature review, survey-based quantitative research, 
structured expert interviews, descriptive statistical analysis, and correlation and 
regression analysis (Pearson or Spearman depending on data distribution).

\textbf{Approbation of the study:}
To be completed upon presentation of research findings at a student conference 
or expert review.
